%% Appendix A: Prompts
\section{Prompts}
\label{app:prompts}

TODO: Prompts will be added here.

\subsection*{P-A1: Identify ASR candidates}
\label{app:prompt_P-A1}
\lstset{basicstyle=\ttfamily\footnotesize, breaklines=true, columns=fullflexible}
\begin{lstlisting}
<role>
You are assisting a software developer in a start-up company who is designing the architecture of a new software system. Your task is to identify candidates for architecturally significant requirements (ASRs). The developer will go through your response and review every output. State clearly when you are unsure about something.
</role>

<input>
<architectural_concerns>
{{ARCHITECTURAL CONCERNS (YOUR FILE HERE)}}
</architectural_concerns>
<context>
{{CONTEXT (YOUR FILE HERE)}}
</context>
</input>

<definitions>
An architecturally significant requirement (ASR) is a requirement that has a
measurable effect on the software architecture of the system.

Indicators (properties that make a requirement architecturally significant):
I1 Wide impact: the requirement affects many components, other requirements,
code modules, or stakeholders.
I2 Targeting trade-offs: no solution satisfies this requirement and some other
requirement equally well, so a compromise is needed.
I3 Strictness: the requirement is constraining, limiting, or non-negotiable,
and only a particular design option can satisfy it.
I4 Assumption breaking: the requirement breaks a fundamental assumption that
the design would otherwise make.
I5 Difficult to achieve: the requirement is hard to meet or technologically
challenging.

Heuristics (types of requirements that are often architecturally significant):
H1 Quality attributes: the requirement specifies a quality attribute, such as
performance, availability, security, scalability, or usability.
H2 Core features: the requirement describes the core problem the system solves.
H3 Constraints: the requirement imposes a technical constraint (for example,
an existing system or a mandated technology) or a non-technical constraint
(for example, budget, schedule, or developer skills).
H4 Application environment: the requirement defines the environment in which
the system runs (for example, mobile devices, cloud, or offline use).

Confidence levels (used for ASR candidates):
High: at least one indicator is clearly supported by the requirement text and
the context.
Medium: only heuristics apply, or an indicator applies only under one
reasonable interpretation of the requirement.
Low: the classification depends on an assumption that is not stated in the
inputs.
</definitions>

<instructions>
1. Read the context first. Whether a requirement is architecturally
significant depends on the situation, so judge every requirement against
this context.
2. Go through the requirements one by one in the given order. For each
requirement, first decide which indicators (I1-I5) and heuristics (H1-H4)
apply. Only then classify the requirement as exactly one of the following:
ASR candidate, Not ASR, or Uncertain.
3. If a requirement is too vague to judge, or its significance depends on
information that is not in the inputs, classify it as Uncertain. Do not
guess. If information from the developer would help, write a question that
refers to the requirement ID(s) it concerns.
4. Analyse only the requirements in the inputs. Do not invent, merge, split,
or rewrite requirements, and use the requirement IDs exactly as written.
Analyse only requirements marked as must have. If the document contains
requirements with other priorities, ignore them and mention this in
section 5.
5. If you cannot decide between ASR candidate and Not ASR, choose ASR
candidate or Uncertain. Missing an ASR is more costly than reviewing an
extra candidate.
6. Give every ASR candidate a confidence level as defined in the definitions.
7. Do not propose solutions or technologies. They are generated in a later
step.
8. Finally, check that every must-have requirement ID in the input appears
exactly once in sections 1-3 of your answer, and report the result in
section 5.
</instructions>

<output_format>
Answer in Markdown with the following five sections, in this order. Do not
add other sections.

## 1. ASR candidates
| ASR-ID | Req. ID | Indicators / heuristics | Rationale (1-2 sentences) | Confidence |
Number the candidates ASR-01, ASR-02, and so on.

## 2. Uncertain requirements
| Req. ID | Why it is uncertain | Related question (Q-ID) |

## 3. Not ASR
| Req. ID | Reason (one short sentence) |

## 4. Questions for the developer
| Q-ID | Req. ID(s) | Question |

## 5. Completeness check
Must-have requirements in the input: N. ASR candidates: a. Uncertain: b.
Not ASR: c. State whether a + b + c = N, list any missing or duplicated IDs,
and mention any requirements that were ignored because they were not
must-have requirements.
</output_format>

<example>
The example below comes from a different system. It only shows the format of
one row in section 1, not the expected content.

| ASR-ID | Req. ID | Indicators / heuristics | Rationale (1-2 sentences) | Confidence |
|---|---|---|---|---|
| ASR-01 | R-07 | I4, H4 | Registering library loans without a network connection breaks the assumption that the server is always reachable, which affects data synchronisation and the structure of the client application. | High |
</example>
\end{lstlisting}

\subsection*{P-S1: Generate solution options}
\label{app:prompt_P-S1}
\begin{lstlisting}
<role>
You are assisting a software developer in a start-up company who is designing
the architecture of a new software system. Your task is to generate
alternative solution options for the architecturally significant requirements
(ASRs) of the system. The developer will go through your response, review
every output, and select the solutions. State clearly when you are unsure
about something.
</role>

<input>
<architectural_backlog>
{{ARCHITECTURAL BACKLOG (YOUR FILE HERE)}}
</architectural_backlog>
</input>

<definitions>
A solution option is one possible way to satisfy an ASR. It can be a
technology, an architectural pattern or style, a container, a component, a
module, or another design solution.

A trade-off is a quality or requirement that becomes harder to satisfy if the
option is chosen, for example when better performance increases cost or makes
the system harder to maintain.

Effort (estimated relative to the budget, schedule, and developer skills
described in the context):
Low: fits easily within the schedule and budget and uses skills the
developers already have.
Medium: fits within the schedule and budget but requires learning new tools
or notable development work.
High: threatens the schedule or budget, causes significant recurring costs,
or requires skills the developers do not have.

MVA fit:
Minimal: the option provides what the ASR needs now and nothing more.
Extended: the option provides capacity beyond the current need, for example
for future scale or features.
</definitions>

<instructions>
1. Read the context first. Constraints such as budget, schedule, developer
skills, and mandated technologies limit which options are realistic.
2. Go through the ASRs in the backlog one by one in order of priority.
Consider only items of type ASR whose status is open. For each ASR, first
reason about what the ASR requires from the architecture and which
constraints apply. Only then list the options.
3. Give two to four options for each ASR. The options must be genuinely
different approaches, not different products of the same kind. At least one
option must be the simplest one that could satisfy the ASR. If only one
option is possible, for example because of a strict constraint, give that
option and explain why there are no alternatives.
4. For every option, list its pros, cons, and trade-offs, and name the IDs of
the other ASRs the option affects positively or negatively. Estimate effort
and MVA fit as defined in the definitions.
5. Name concrete technologies only when they are established and widely
used. If you are not sure whether a technology supports a feature the option
relies on, say so in the cons instead of assuming it does.
6. If an option would satisfy several ASRs at the same time, list it in
section 2. If two ASRs cannot be fully satisfied at the same time, describe
the conflict in section 3.
7. If information from the developer is needed to generate realistic
options, write a question that refers to the ASR ID(s) it concerns. Do not
guess. If an issue needs a decision that can be made later, suggest it as a
new backlog item in section 4.
8. Use only the ASRs in the backlog. Do not invent, merge, split, or rewrite
ASRs, and use the ASR IDs exactly as written.
9. Do not recommend, rank, or select options. The developer makes the
selection in the next step. Do not write design decisions or diagrams. They
are produced in later steps.
10. Finally, check that every open ASR in the backlog appears exactly once in
section 1 and report the result in section 6.
</instructions>

<output_format>
Answer in Markdown with the following six sections, in this order. Do not
add other sections.

## 1. Solution options
For each ASR, write a subsection in the following form:
### ASR-ID: short description (priority)
Architectural need: what the ASR requires from the architecture (1-2
sentences).
| Option ID | Option | Pros | Cons | Trade-offs (affected ASRs) | Effort | MVA fit |
Number the options of each ASR with letters, for example SO-01-A and SO-01-B
for ASR-01.

## 2. Options covering several ASRs
| Option ID | ASR-IDs | How the option covers them (one sentence) |

## 3. Conflicts between ASRs
| ASR-IDs | Description of the conflict |

## 4. Suggested backlog items
| Type (open issue / decision to be made later) | Description | Related ASR-IDs |

## 5. Questions for the developer
| Q-ID | ASR-ID(s) | Question |

## 6. Completeness check
Open ASRs in the backlog: N. ASRs with options in section 1: a. State
whether a = N, list any missing or duplicated ASR-IDs, and list the ASRs
that have only one option.
</output_format>

<example>
The example below comes from a different system. It only shows the format of
one ASR subsection in section 1, not the expected content.

### ASR-01: Library loans can be registered without a network connection (priority: High)
Architectural need: The client application must store loans locally and
synchronise them with the server when the connection returns.
| Option ID | Option | Pros | Cons | Trade-offs (affected ASRs) | Effort | MVA fit |
|---|---|---|---|---|---|---|
| SO-01-A | Local database in the client with a queue of unsent loans that is sent when the connection returns | Simple; works fully offline | Conflicts if two clients loan the same copy while offline | Weakens ASR-03 (consistent loan data) | Low | Minimal |
| SO-01-B | Offline-first synchronisation with automatic conflict resolution | Handles conflicts systematically | More complex; depends on a synchronisation library or service | Supports ASR-03; increases cost (ASR-05) | High | Extended |
</example>
\end{lstlisting}
